\documentclass[12pt]{article}
\usepackage[a3paper, margin=1in]{geometry}
\usepackage{amsmath, amssymb, amsthm, ulem}

\begin{document}

\section*{Domácí úkol 5.1}

Vypracoval: Daniel \textit{"Randál"} Ransdorf \hfill Podpis: \rule{4cm}{0.4pt}

\begin{enumerate}
  \item
  
    \begin{enumerate}
      \item \textbf{Hlavní idea / pozorování}: Budeme hledat zprava inverzní matici.
        Všimněme si, že když vezmeme libovolný sloupcový vektor kromě prvního a odečteme
        sloupcový vektor hned vedle nalevo, získáme aritmetický vektor s nulama a jedním číslem
        uprostřed.
        Sloupcovým pohledem si tedy všimněme, že i-tý sloupec jednotkové matice jde složit
        lineární kombinací "i-tý sloupec zprava" - "sloupec vedle vlevo" / i, kromě
        posledního sloupce, který se dá vytvořit "první sloupec" / n.
        A tedy inverzní matici tedy zkonstruujeme ze sloupcových vektorů
        ve tvaru $\frac{1}{i}(0,...,0,-1,1,0,...,0)^T$.
        
      \item \textbf{Řešení}:

        Označme matici v zadání $A$ zjevně typu $n \times n$,
        aritmetický vektor i-tého sloupce matice $A$ označme $\vec{a}_i$. Hledejme matici
        $X$ typu $n \times n$, aritmetický vektor i-tého sloupce matice $X$ označme $\vec{x}_i$.
        V $\vec{a}_i$ je ve "spodní" části řada čísel $n-i+1, n-i+2, ..., n$
        a v "horní" části příslušný počet nul pro doplnění do rozměru matice.
        Ve sloupcovém vektoru $\vec{a}_i$ je ve "spodní" řadě právě $n-(n-i+1)+1=i$ čísel, tím pádem bude v "horní" části
        $n-i$ nul.
        
        Označme $O_x$ nulový aritmetický vektor o $x$ složkách, pro záporné $x$ bude $O_x$ reprezentovat vektor o 0 složkách. 
        Dle úvahy určeme předpis sloupcového vektoru matice A:

        \[
          \vec{a}_i = 
          \left(
            \begin{array}{c}
              \left.\begin{array}{c}
              0 \\ 
              \vdots \\ 
              0
              \end{array}\right\} \! n-i \\
              n-i+1 \\ n-i+2 \\ \vdots \\ n
            \end{array}
          \right)=
          \left(\begin{array}{c} 
            O_{n-i} \\ \hline n-i+1 \\ n-i+2 \\ \vdots \\ n
          \end{array}\right)
        \]

        Vytvořme si rovnou vzorec pro odečtení dvou vektorů vedle sebe pro $i \geq 2$:

        \begin{align*}
          \vec{a}_{i} - \vec{a}_{i-1} &= 
            \left(\begin{array}{c} 
              O_{n-i} \\ \hline n-i+1 \\ n-i+2 \\ \vdots \\ n
            \end{array}\right)
            -
            \left(\begin{array}{c} 
              O_{n-i+1} \\ \hline n-i+2 \\ n-i+3 \\ \vdots \\ n
            \end{array}\right)
          =
            \left(\begin{array}{c} 
              O_{n-i} \\ \hline n-i+1 \\ n-i+2 \\ \vdots \\ n
            \end{array}\right)
            -
            \left(\begin{array}{c} 
              O_{n-i} \\ \hline 0 \\ n-i+2 \\ n-i+3 \\ \vdots \\ n
            \end{array}\right)
          =
            \left(\begin{array}{c} 
              O_{n-i} - O_{n-i} \\ \hline (n-i+1 )- 0 \\ (n-i+2) - (n-i+2) \\ (n-i+3) - (n-i+3) \\ \vdots \\ n-n
            \end{array}\right)
          \\
          \vec{a}_{i} - \vec{a}_{i-1} &= 
            \left(\begin{array}{c} 
              O_{n-i} \\ \hline n-i+1 \\ \hline O_{n-(n-i+2)+1}
            \end{array}\right)
          =
            \left(\begin{array}{c} 
              O_{n-i} \\ \hline n-i+1 \\ \hline O_{i-1}
            \end{array}\right) \quad (*)
        \end{align*}

        Podobným způsobem si označme aritmetický vektor i-tého sloupce jednotkové matice typu $n \times n$ jako $\vec{j_i}$:

        \[
          \vec{j_i} = \left(\begin{array}{c} 
              O_{i-1} \\ \hline 1 \\ \hline O_{n-i}
            \end{array}\right)
        \]

        Hledaná matice $X$ musí splňovat následující vztah, který dále rozepišme sloupcovým pohledem násobení matic:

        \begin{align*}
          AX &= I_n \\
          \Leftrightarrow A\vec{x}_i &= \vec{j}_i
        \end{align*}

        Zapišme předpis pro $\vec{x}_i$ dle pozorování.

        \begin{align*}
          \vec{x}_i  = 
            \begin{cases}
            i < n: & \dfrac{1}{i} \cdot \left(\begin{array}{c} 
                O_{n-i-1} \\ \hline -1 \\ 1 \\ \hline O_{i-1}
              \end{array}\right) 
            \\[6pt]
            i=n: & \left(\begin{array}{c} 
                \frac{1}{n} \\ \hline O_{n-1}
              \end{array}\right) 
            \end{cases}
        \end{align*}

        Ukažme, že součin $A\vec{x}_i$ je roven $\vec{j}_i$:
          \begin{enumerate}
            \item $i<n$:
              
              \begin{align*}
                A\vec{x}_i &= \underbrace{O_n + ... + O_n}_{n-i-1} \;\;+\;\;  \frac{1}{i} \cdot (-1) \cdot \vec{a}_{n-i} \;\;+\;\; \frac{1}{i} \cdot \vec{a}_{n-i+1}
                \;\;+\;\; \underbrace{O_n + ... + O_n}_{i-1} \\
                &= \frac{1}{i} (\vec{a}_{n-i+1} - \vec{a}_{n-i})
                \overset{(*)}{=} \frac{1}{i} \left(\begin{array}{c} 
                  O_{n-(n-i+1)} \\ \hline n-(n-i+1)+1 \\ \hline O_{(n-i+1)-1}
                \end{array}\right)
                = \frac{1}{i} \left(\begin{array}{c} 
                  O_{i-1} \\ \hline i \\ \hline O_{n-i}
                \end{array}\right) \\
                A\vec{x}_i &= \left(\begin{array}{c} 
                  O_{i-1} \\ \hline 1 \\ \hline O_{n-i}
                \end{array}\right) = \vec{j}_i
              \end{align*}
            
            \item $i=n$:
            
              \begin{align*}
                A\vec{x}_n &= \frac{1}{n} \cdot \vec{a}_1 \;\;+\;\;  \underbrace{O_n + ... + O_n}_{n-i} \\
                A\vec{x}_n &= \frac{1}{n} \cdot \vec{a}_1 = \frac{1}{n} \cdot \left(\begin{array}{c} 
                  O_{n-1} \\ \hline n
                \end{array}\right) = \left(\begin{array}{c} 
                  O_{n-1} \\ \hline 1
                \end{array}\right) = \vec{j}_n
              \end{align*}
          \end{enumerate}

        Našli jsme tedy předpis sloupců matice $X$, který zajistí, že $AX=I_n$.
        Tím pádem je matice $A$ invertovatelná. Jelikož je $A$ čtvercová a invertovatelná matice,
        pak $X$ je inverzní matice k $A$ zleva i zprava.
        Uveďme tedy varianty zápisu:
        \renewcommand{\arraystretch}{1.4}

        \begin{align*}
          X &= \left(\begin{array}{ccc} 
              \vec{x}_1 \;| & ... & |\; \vec{x}_n
            \end{array}\right) \\
          X &= \left(
            \begin{array}{ccccccc} 
              0&0&0& \cdots &0 &\frac{-1}{n-1}&\frac{1}{n} \\
              0&0&0& \cdots & \frac{-1}{n-2} &\frac{1}{n-1}&0 \\
              0&0&0& \cdots & \frac{1}{n-2} &0&0 \\
              \vdots & \vdots & \vdots & \cdots & \vdots & \vdots & \vdots \\
              0&0& \frac{-1}{3} & \cdots & 0&0&0 \\
              0&\frac{-1}{2}& \frac{1}{3} & \cdots & 0&0&0 \\
              -1&\frac{1}{2}& 0 & \cdots & 0&0&0 \\
              1&0& 0 & \cdots & 0&0&0 \\
            \end{array}
          \right)
        \end{align*}

          
    \end{enumerate}

\end{enumerate}

\end{document}
